\documentclass[10pt,a4paper]{amsart}
\usepackage[margin=19mm]{geometry}
\usepackage{amsmath,amssymb,mathtools}
\usepackage[scr=rsfs,cal=euler]{mathalfa}
\usepackage{xfrac}
\usepackage[unicode,hidelinks,bookmarks=false]{hyperref}
\newcommand{\ie}{i.e.\ }
\newcommand{\cf}{cf.\ }
\newcommand{\resp}{respectively\ }
\newcommand{\Subsection}{\subsection}
\providecommand{\nobreakdash}{\nobreak\hbox{-}}
\setlength{\emergencystretch}{2em}
\multlinegap0pt
\usepackage{bm}
\usepackage{bbm}
\usepackage{dsfont}
\usepackage{xspace}
\usepackage{cancel}
\usepackage{mathtools}
\usepackage{amsthm}
\makeatletter
\@ifundefined{lemma}{\newtheorem{lemma}{Lemma}}{}
\@ifundefined{theorem}{\newtheorem{theorem}{Theorem}}{}
\makeatother
\renewcommand{\epsilon}{\varepsilon}
\title[Higher regularity of solutions of an iterative functional eq]{Higher regularity of solutions of an iterative functional equation}
\author{Liang Feng, Xiao Tang}
\date{Source: arXiv:2604.14244v1 (CC BY 4.0)}
\begin{document}
\maketitle

\noindent\emph{Source and licence.} Higher regularity of solutions of an iterative functional equation. Version arXiv:2604.14244v1 (CC BY 4.0), distributed under the Creative Commons Attribution 4.0 International licence, as recorded in the arXiv licence metadata for this exact version and shown on the abstract page https://arxiv.org/abs/2604.14244v1. Full rights notice: licenses/arxiv-2604.14244-CC-BY-4.0.txt.\par\smallskip

\noindent\textit{Editorial scope.} Counted unit: the Lemma whose source label is \texttt{coninuous of fg} in the article (source0.tex lines 262--270, source label \texttt{coninuous of fg}), delivered together with the article's own complete original proof of it (source0.tex lines 271--330, one \texttt{proof} environment of 60 lines). No mathematics is written, repaired or re-derived: statement and proof are the article's own text line for line. Required context retained uncounted: o4 (lines 134--136); th (lines 143--161). Every definition, display and label of those lines is retained unchanged. Omitted from this package and not redistributed: 1--133, 137--142, 162--261, 331--548. Nothing inside a retained excerpt is omitted or altered: every retained body is delivered exactly as the article's lines are written, so no omission receipt of any kind accompanies this package. Citations: the retained excerpts carry no citation command at all, so no bibliography entry of the article is transcribed and no reference is redistributed. Reference adaptation: none was needed and none was made; every reference inside the delivered unit resolves inside it, so no unresolved reference or citation is produced. Printed slips kept literal: no printed word, symbol, hypothesis or number of the counted statement or of its proof was altered. Layout and class: the article's own class is \texttt{amsart} with packages including amssymb, amsmath, euscript, hyperref, graphicx, ulem, chngpage, xcolor, enumitem; the wrapper is the lane's pinned \texttt{amsart} editorial wrapper at 10pt on A4, which restates the theorem-like environments the retained text needs and adds the article's own macro definitions that the retained text needs (\texttt{\textbackslash epsilon}), with extra wrapper packages (bm, bbm, dsfont, xspace, cancel, mathtools). The delivered numbering is the unit's own. Assets: no retained span contains a figure environment, a graphics-inclusion command, a PDF-page inclusion, a table or a TikZ picture; no image file is copied into this package, the package declares no asset dependency and no image file is redistributed with it. Licence: version arXiv:2604.14244v1 of this article is distributed under the Creative Commons Attribution 4.0 International licence, as recorded in the arXiv licence metadata for this exact version and shown on the abstract page https://arxiv.org/abs/2604.14244v1. A full rights notice is delivered at licenses/arxiv-2604.14244-CC-BY-4.0.txt. Characters: every retained excerpt is pure ASCII, so the delivered engine, XeLaTeX with its default Latin Modern text font, reports no missing character.
\begin{equation}\label{o4}
    \phi^2(x) = h(\phi(f(x))) +g(x), \qquad x\in \mathbb{R},
\end{equation}

\begin{theorem}\label{th}
Assume that functions $h:\mathbb{R} \to \mathbb{R}, f:\mathbb{R} \to  \mathbb{R}$ and $g:\mathbb{R} \to \mathbb{R}$ are of class $C^{n}$ such that
\begin{equation}\label{o}
\inf_{x\in \mathbb{R}}|h'(x)|\ge m, 
~\inf_{x\in \mathbb{R}}|f'(x)|\ge \alpha,
\end{equation}
\begin{equation}\label{gjie}
\max\Big\{\sup_{x\in \mathbb{R}}|g'(x)|,\sup_{x\in \mathbb{R}}|g''(x)|,...,\sup_{x\in \mathbb{R}}|g^{(n)}(x)| \Big\}~\le \beta, ~\sup_{x\in \mathbb{R}}|g(x)|<+\infty,~
\end{equation}
 where $m > 1$, $\alpha > 0$ are real constants, and $\beta > 0$ is a sufficiently small real constant.
Furthermore, 
\begin{equation}\label{oo}
\max_{1\le k\le n}\Big\{ \sup_{x\in \mathbb{R}}|h^{(k)}(x)|,
  \sup_{x\in \mathbb{R}}|f^{(k)}(x)|
 \Big\}< +\infty.
 \end{equation}
Then the functional equation \eqref{o4}
has a bounded $C^{n}$ solution whose the first up to $n$-th derivatives are bounded.
\end{theorem}

\begin{lemma} \label{coninuous of fg}
Under the assumption of Theorem \ref{th}, for every $j=1,2,\cdots,n,$ the maps $\mathcal{F}_j:  C^{0,1}_b(\mathbb{R};L) \to C^0_b(\mathbb{R})$ defined by \[\mathcal{F}_j(\phi):=(h^{-1})^{(j)}\circ (\phi^{2} -g )\circ f^{-1} \qquad \text{for every } \phi\in C^{0,1}_b(\mathbb{R};L),  \]
where $(h^{-1})^{(j)}$ denotes the $j$-th derivative, the map $\mathcal{G}: C^0_{b,\rho }(\mathbb{R})  \times C^{0,1}_b(\mathbb{R};L) \to  C^0_{b,\rho }(\mathbb{R})$ defined by 
\[\mathcal{G} (\psi,\phi):= \psi \circ \phi \circ f^{-1} \qquad  \]
for every $(\psi,\phi) \in C^0_{b,\rho }(\mathbb{R})\times C^{0,1}_b(\mathbb{R};L)$, and the map $\mathcal{B}:  C^0_{b,\rho }(\mathbb{R}) \to  C^0_{b,\rho }(\mathbb{R})$ defined by
$$ \mathcal{B}(\phi) := \phi \circ f^{-1} \qquad\text{ for every } \phi \in C_{b,\rho}^{0}(\mathbb{R}), $$
are all continuous.

\end{lemma}

\begin{proof}
 First, we prove that $\mathcal{F}_j: C^{0,1}_b(\mathbb{R};L) \to C^0_b(\mathbb{R})$ is continuous for each $j=1,2,\cdots,n$. For any $\phi_0 \in C_{b}^{0,1}(\mathbb{R};L)$, we need to prove that, for any $\epsilon >$ 0, there exists $\delta >0$, such that 
  $$\|(h^{-1})^{(j)} \circ (\phi^{2} \circ f^{-1}-g \circ f^{-1})-(h^{-1})^{(j)} \circ (\phi_0^{2} \circ f^{-1}-g \circ f^{-1})\|<\epsilon$$ whenever $\|\phi-\phi_0\| < \delta$.
  Since the functions $(h^{-1})^{(j)}$ is continuous on $\mathbb{R}$, it is uniformly continuous on the bounded closed interval $I:= [-\|\phi_0\|-\|g\|-1,\|\phi_0\|+\|g\|+1]$. Thus, there exists a $\delta_0  \in (0,1)$ such that
  \begin{equation}\label{unifrom-conti-h}
   |(h^{-1})^{(j)}(x_1)-(h^{-1})^{(j)}(x_2)|< \frac{\epsilon}{2} \text{ whenever } |x_1-x_2|< \delta_0 \text{ and } x_1,x_2 \in I.   
  \end{equation}
  When $\|\phi-\phi_0\| < \frac{\delta_0}{L+1},$ we see that
  \begin{align}\label{pp}
|\phi^{2} \circ f^{-1}(x)-g \circ f^{-1}(x)| &\le \sup_{x \in \mathbb{R}}|\phi^{2} \circ f^{-1}(x)-g \circ f^{-1}(x)|\nonumber
\\
&=
\|\phi^{2} \circ f^{-1}-g \circ f^{-1}\| \nonumber\\
&\le \|\phi\| +\|g\| \le \|\phi-\phi_0\| +\|\phi_0\|+\|g\|
\nonumber\\
& \le \frac{\delta_0}{L+1} +\| \phi_0 \|+\|g\|
\nonumber\\
&\le 1+\| \phi_0 \|+\|g\| \qquad  \text{ for all } x \in \mathbb{R},
  \end{align}
that is, $ \phi^{2} \circ f^{-1}(x)-g \circ f^{-1}(x) \in I$ for all $ x \in \mathbb{R}$; and that
   \begin{align}\label{mm}
  |\phi^{2} \circ f^{-1}(x)- \phi_0^{2} \circ f^{-1}(x)| &\le \|\phi^{2}-\phi_0^{2}\| \le \|\phi^{2}-\phi \circ \phi_0\| +\|\phi \circ \phi_0-\phi_0^{2}\| \nonumber\\ &\le (L+1)\|\phi-\phi_0\| < \delta_0 \qquad  \text{ for all } x \in \mathbb{R}.
  \end{align}
  Consequently,  when $\|\phi-\phi_0\| < \frac{\delta_0}{L+1},$ by \eqref{unifrom-conti-h}, \eqref{pp} and \eqref{mm}, 
  \begin{align*}
  &\|(h^{-1})^{(j)}\circ (\phi^{2} \circ f^{-1}-g \circ f^{-1})-(h^{-1})^{(j)}\circ (\phi_0^{2}\circ f^{-1}-g \circ f^{-1})\| \\
  &= \sup_{x \in \mathbb{R}}|(h^{-1})^{(j)}\circ (\phi^{2} \circ f^{-1}(x)-g \circ f^{-1}(x))-(h^{-1})^{(j)}\circ (\phi_0^{2} \circ f^{-1}(x)-g \circ f^{-1}(x))|\\
  &\le \frac{\epsilon}{2} < \epsilon.
  \end{align*}
  The continuity of $\mathcal{F}_j(\phi), j=1,2,\cdot \cdot \cdot,n,$ is proved.  
  Next, we prove that $\mathcal{G} (\psi,\phi)$ is continuous. For any $(\tilde{\psi},\tilde{\phi}) \in C_{b,\rho}^{0} (\mathbb{R}) \times C_{b}^{0,1} (\mathbb{R};L)$, we need to prove that, for any $\epsilon >$ 0, there exists $\delta >0$, such that
  $$\|\psi \circ \phi \circ f^{-1}-\tilde{\psi} \circ \tilde{\phi} \circ f^{-1}\| < \epsilon \text{ whenever } d((\psi,\phi),(\tilde{\psi},\tilde{\phi})):=\max\{\| \phi-\tilde \phi\|,\| \psi-\tilde \psi\|\}<\delta.$$
Since 
\begin{align}
    &\|\psi \circ \phi \circ f^{-1}-\tilde{\psi} \circ \tilde{\phi} \circ f^{-1}\|\notag\\ 
    &\le \|\psi \circ \phi \circ f^{-1}-\tilde{\psi} \circ \phi \circ f^{-1}\|+\|\tilde{\psi} \circ \phi \circ f^{-1}-\tilde{\psi} \circ \tilde{\phi} \circ f^{-1}\|\notag\\
    &\le \|\psi-\tilde \psi\|+\|\tilde{\psi} \circ \phi \circ f^{-1}-\tilde{\psi} \circ \tilde{\phi} \circ f^{-1}\|. \notag
\end{align}
Thus prove the continuity of $\mathcal{G}$ is complete under $\tilde{\psi} \circ \phi \circ f^{-1}$ is continuous. Next, we just need to prove $\tilde{\psi} \circ \phi \circ f^{-1}$ is continuous. In other words, for every $\epsilon >0$, there exists $\delta >0,$ such that
$$\|\tilde{\psi} \circ \phi \circ f^{-1}-\tilde{\psi} \circ \tilde{\phi} \circ f^{-1}\| < \epsilon \text{ whenever } \|\phi-\tilde \phi\|<\delta.$$
Since the function $\tilde \psi$ is continuous on $\mathbb{R}$, it is uniformly continuous on bounded closed interval $I: =[-\|\tilde{\phi}\|-1,\|\tilde{\phi}\|+1]$. That is to say, for any $\epsilon >0$, there exists $0<\delta_0 <1$, such that
\begin{equation}\label{A}
    |\tilde \psi(x_1)-\tilde \psi(x_2)| < \frac{\epsilon}{2} \text{ whenever } |x_1-x_2|<\delta_0  \text{ and } x_1,x_2 \in I.\\
\end{equation}
In that case, when $\|\phi-\tilde{\phi}\| < \delta_0$,
\begin{equation}\label{AA}
    |\phi \circ f^{-1}(x)| \le \|\phi\| \le \|\phi-\tilde{\phi}\|+\|\tilde{\phi}\| < 1+ \|\tilde{\phi}\|
\qquad \text{ for all } x \in \mathbb{R},
\end{equation}
that is, $\phi \circ f^{-1}(x) \in I$  \qquad \text{ for all } $x \in \mathbb{R};$
and that 
\begin{equation}\label{AAA}
    |\phi \circ f^{-1}(x)-\tilde{\phi} \circ f^{-1}(x)| \le \|\phi-\tilde{\phi}\| < \delta_0 \qquad \text{ for all } x \in \mathbb{R}.
\end{equation}
Thus, when $\|\phi-\tilde{\phi}\| < \delta_0, $ by \eqref{A},\eqref{AA} and \eqref{AAA}, 

$\|\tilde \psi \circ \phi \circ f^{-1}-\tilde \psi \circ \tilde{\phi} \circ f^{-1}\| =\sup_{x \in \mathbb{R}}|\tilde\psi \circ \phi \circ f^{-1}(x)-\tilde\psi \circ \tilde{\phi} \circ f^{-1}(x)| \le \frac{\epsilon}{2} < \epsilon.$
The continuity of $\mathcal{G} (\psi,\phi)$ is proved. 
Similarly, we are able to show that $\mathcal{B}$ is also continuous with respect to $\phi$. The proof is complete.
\end{proof}

\end{document}
